\documentclass[11pt]{article}
\usepackage[utf8]{inputenc}
\usepackage[T1]{fontenc}
\usepackage{babel}
\usepackage{lmodern}
\usepackage{amsmath,amssymb}
\usepackage{hyperref}
\usepackage{mathptmx}

% THEOREM CONFIGS (ENGLISH)
\usepackage[thmmarks]{ntheorem}

\theoremstyle{plain}
\theoremheaderfont{\normalfont\bfseries}
\theorembodyfont{\normalfont}
\theoremseparator{:}
\theorempreskip{\topsep}
\theorempostskip{\topsep}

\theoremsymbol{\ensuremath{\bigcirc}}
\theorembodyfont{\normalfont}
\newtheorem{definition}{Definition}[section]

\theoremsymbol{\ensuremath{\triangle}}
\theorembodyfont{\normalfont}
\newtheorem{note}[definition]{Note}

\theoremsymbol{}
\theorembodyfont{\normalfont}
\newtheorem{proposition}[definition]{Proposition}

\theoremsymbol{\ensuremath{\square}}
\theorembodyfont{\normalfont}
\newtheorem{freeprop}[definition]{Proposition}

\theoremstyle{nonumberplain}
\theoremheaderfont{\normalfont\bfseries}
\theorembodyfont{\normalfont}
\theoremsymbol{\ensuremath{\blacksquare}}
\newtheorem{proof}[definition]{Proof}

% Generic Macros (english)
\newcommand{\naturals}{\mathbb{N}}
\newcommand{\integers}{\mathbb{Z}}
\newcommand{\rationals}{\mathbb{Q}}
\newcommand{\reals}{\mathbb{R}}
\newcommand{\complexs}{\mathbb{C}}
\newcommand{\set}[1]{\{#1\}}
\newcommand{\inv}[1]{{#1}^{-1}}

\title{Notes for Mathematical Foundations of Automata Theory by Jean-Éric Pin}
\author{Gian Carlo}
\date{\today}

\begin{document}
\maketitle

\section{Sets, Functions, and Relations}

\subsection{Functions}

Usual definitions of functions and relations are given. Some key properties of functions, with respect to composition, are listed:

\begin{enumerate}
    \item Injective functions are left cancellative;
    \item Surjective functions are right cancellative; and
    \item Bijective functions are invertible.
\end{enumerate}

Another property is stated here fully (except for variable names):

\begin{freeprop}
    Let $p : E \to F$ and $q : F \to G$ be two mappings and let $f = q \circ p$ be their composite. Then,
    \begin{enumerate}
        \item If $f$ is injective, then $p$ is injective. Furthermore, $p$ being surjective implies $q$ being injective; and
        \item If $f$ is surjective, then $q$ is surjective. Furthermore, $q$ being injective implies $p$ is surjective.
    \end{enumerate}
\end{freeprop}

Another useful result regarding functions is stated (and here repeated verbatim):

\begin{freeprop}
    Let $E$ and $F$ be two finite sets such that $|E| = |F|$ and let $\phi : E \to F$ be a function. The following conditions are equivalent: \begin{enumerate}
        \item $\phi$ is injective,
        \item $\phi$ is surjective,
        \item $\phi$ is bijective.
    \end{enumerate}
\end{freeprop}

\subsection{Relations}

The notions of injectivity and surjectivity are extended into relations in the natural way, mirroring their function form definitions. Some equivalent definitions are given as propositions. For that, we repeat the definition of partial function in the book: a function such that each element has \textbf{at most} one image (which means it could have none).

\begin{freeprop}
    A relation is injective if and only if its inverse is a partial function. (In other words, a partial function induces a relation from the codomain into the domain.)
\end{freeprop}

\begin{freeprop}
    Let $\tau : E \to F$ be a relation. Then $\tau$ is injective if and only if, for all $X, Y \subseteq E$ we have that $X \cap Y = \emptyset$ implies $\tau(X) \cap \tau(Y) = \emptyset$.
\end{freeprop}

This next proposition is a bit more involved. For this, let $e(x) = x$ be the identity function. Specifically, $e_F$ is the identity restricted to the set $F$.

\begin{freeprop}
    Let $\tau : E \to F$ be a relation. The following conditions are equivalent: \begin{enumerate}
        \item $\tau$ is injective,
        \item $\inv \tau \circ \tau = e_{Dom(\tau)}$,
        \item $\inv \tau \circ \tau \subseteq e_E$.
    \end{enumerate}
\end{freeprop}

There are two useful consequences of this proposition:

\begin{freeprop}
    Let $\tau : E \to F$ be a relation. The following conditions are equivalent: \begin{enumerate}
        \item $\tau$ is a partial function,
        \item $\tau \circ \inv \tau = e_{Im(\tau)}$,
        \item $\tau \circ \inv \tau \subseteq e_F$.
    \end{enumerate}
\end{freeprop}

\begin{freeprop}
    Let $\tau : E \to F$ be a relation. Then $\tau$ is a surjective partial function if and only if $\tau \circ \inv \tau = e_F$.
\end{freeprop}

This last proposition is used in the following form:

\begin{freeprop}
    Let $p : F \to G$ and $q : E \to F$ be two functions and let $f = p \circ q$. If $q$ is surjective, the relation $f \circ \inv q$ is equal to $p$.
\end{freeprop}

Composition of same-jectivity relations yield again a same-jective relation. We can adapt a similar proposition stated for functions before:

\begin{freeprop}
    Let $p : E \to F$ and $q : F \to G$ be two relations and let $f = q \circ p$ be their composition.
    \begin{enumerate}
        \item If $f$ is injective and $\inv q$ surjective, then $p$ is injective. Furthermore, if $p$ is a surjective partial function, then $q$ is injective.
        \item If $f$ is surjective, then $\inv q$ is surjective. Furthermore, if $q$ is injective of $F$, then $q$ is surjective.
    \end{enumerate}
\end{freeprop}

This next proposition is given as-is.

\begin{freeprop}
    Let $E, F,$ and $G$ be three sets and $p : G \to E$ and $q : G \to F$ be two functions. Suppose $p$ is surjective and that for every $s, t \in G$ we have that $p(s) = p(t)$ implies $q(s) = q(t)$. Then the relation $q \circ \inv p : E \to F$ is a function. 
\end{freeprop}

Finally, we show identities regarding using relations together with set operations.

\begin{freeprop}
    Let $t : E \to F$ be an injective relation. Then for every $X, Y \subseteq E$ we have: \begin{enumerate}
        \item $t(X \cup Y) = t(X) \cup t(Y)$,
        \item $t(X \cap Y) = t(X) \cap t(Y)$, and
        \item $t(X - Y) = t(X) - t(Y)$.
    \end{enumerate}
\end{freeprop}

In particular, the equation $t(X \cup Y) = t(X) \cup t(Y)$ holds for relations in general, not just injective.

\begin{freeprop}
    Let $f : E \to F$ be a function. Then, for every $X, Y \in \subseteq E$, the following holds: \begin{enumerate}
        \item $\inv f (X \cup Y) = \inv f (X) \cup \inv f (Y)$, 
        \item $\inv f (X \cap Y) = \inv f (X) \cap \inv f (Y)$, and
        \item $\inv f (X^c) = \inv f (X)^c$.
    \end{enumerate}
\end{freeprop}

This subsection ends on the following proposition:

\begin{freeprop}
    Let $f : E \to F$ be a partial function. Then for every $X \subseteq E$ and $Y \subseteq F$, we have that \begin{equation*}
        f(X) \cap Y = f(X \cap \inv f (Y))
    \end{equation*}
\end{freeprop}

\subsection{Ordered sets}

A relation $E \to E$ is a \textbf{preorder} if it is reflexive and transitive, an \textbf{order} if it is an antisymmetric preorder, and an \textbf{equivalence relation} if it is a symmetric preorder.

We will order relations by inclusion. If $P$ and $Q$ are two relations on some set $S$, we say that $P$ \textbf{refines} $Q$ (or that $P$ \textbf{is thinner} than $Q$, or that $Q$ is \textbf{coarser} than $P$) if and only if, for each $s, t \in S$, we have that $s P t$ implies $s Q t$. Note that equality is the thinnest preorde and that the universal relation (complete graph) is the coarsest.

\begin{freeprop}
    Any intersection of preorders is a preorder. Any intersection of equivalence relations is an equivalence relation.
\end{freeprop}

Given a set $R$ of relations on a set $E$, there is a \textbf{least} preorder (equiv. relation) containing all the relations on $E$. We will call this the preorder (equiv. relation) generated by $R$.

\begin{freeprop}
    Let $P$ and $Q$ be two preorders [equivalence relations] on a set $E$. If they commute, the preorder [equivalence relation] generated by $P$ and $Q$ is equal to $P \circ Q$.
\end{freeprop}

Here, commuting means that $x (P \circ Q) y$ happens if and only if $x (Q \circ P) y$.

\subsection{Exercises}

Exercise 1 is repeated application of proposition 1.1 together with associating composition.



\end{document}